\documentclass[11pt,letterpaper]{article}

\usepackage[utf8]{inputenc}
\usepackage{geometry}
\geometry{letterpaper, margin=0.75in, top=0.7in, bottom=0.7in}
\usepackage{titlesec}
\usepackage{enumitem}
\usepackage{amsmath,amssymb}
\usepackage{tcolorbox}
\usepackage[hidelinks]{hyperref}
\usepackage{fancyhdr}
\usepackage{microtype}

\titleformat{\section}{\large\bfseries\scshape}{}{0em}{}[\vspace{-2pt}\titlerule\vspace{3pt}]
\titleformat{\subsection}{\normalsize\bfseries}{}{0em}{}

\pagestyle{fancy}
\fancyhf{}
\renewcommand{\headrulewidth}{0pt}
\cfoot{\footnotesize \textcolor{gray}{CloudFlip Researcher Notes \& Interview Defense Guide -- Bhavya Sri Boddeti}}

\tcbset{
  colback=gray!5!white,
  colframe=gray!75!black,
  fonttitle=\bfseries,
  arc=3mm
}

\begin{document}

\begin{center}
  {\LARGE \textbf{CloudFlip: Researcher Notes \& Mathematical Defense Guide}} \\[4pt]
  {\large \textbf{Bhavya Sri Boddeti}} \\[2pt]
  {\small St. Ann's College for Women (Autonomous) \quad$|$\quad Department of Computer Science} \\[2pt]
  {\small Research Preprint: DOI \href{https://doi.org/10.5281/zenodo.23136467}{10.5281/zenodo.23136467} \quad$|$\quad IEEE CCGrid 2027 Submission \#39}
\end{center}

\vspace{2pt}
\hrule height 0.8pt
\vspace{8pt}

\section{The Core Problem in Plain English}
\textbf{The Real-World Bottleneck:} When a cloud application experiences a sudden surge (e.g., Black Friday, flash news, ticketing):
\begin{enumerate}[itemsep=1pt]
  \item \textbf{Metric Telemetry Delay (1--3 minutes):} AWS CloudWatch does not update instantly. It aggregates metrics over 1-to-3-minute periods.
  \item \textbf{Instance Cold-Start Delay (3--5 minutes):} Even after AWS triggers an Auto Scaling Group (ASG) scale-out, launching a new EC2 virtual machine or container takes 3 to 5 minutes (OS boot, pulling Docker images, running entrypoint scripts, and passing ELB health checks).
  \item \textbf{The Saturation Gap (5--8 minutes):} During these unmitigated 5 to 8 minutes, the existing cluster is completely overwhelmed. Server queues overflow, requests time out (HTTP 504), and 99th-percentile (p99) tail latency explodes from 180~ms to over 4,200~ms.
  \item \textbf{Industry Workaround (Static Over-Provisioning):} Because reactive scaling is too late, companies keep 20--50\% surplus idle servers running 24/7 just in case. On a 20-node cluster, this wastes roughly \$670/month in idle compute.
\end{enumerate}

\section{The Solution: How CloudFlip Works}
CloudFlip replaces reactive target tracking with \textbf{predictive pre-warming paired with closed-loop remediation}:
\begin{itemize}[itemsep=1pt]
  \item It continuously monitors incoming traffic trends.
  \item It calculates the rate of acceleration (velocity slope).
  \item It predicts capacity needs \textbf{15 to 30 minutes in advance}.
  \item It orders new instances ahead of time so they are \textbf{already warm and running} when the surge arrives.
  \item It monitors live latency: if p99 latency approaches 180~ms, it enforces an emergency step-up.
\end{itemize}

\section{The Mathematics Explained Simply}

\subsection{Formula 1: Exponential Moving Average (EMA) --- Smoothing the Noise}
\begin{equation}
S_t = \alpha x_t + (1 - \alpha) S_{t-1}
\end{equation}
\begin{tcolorbox}[title=Intuition \& Why We Use It]
\textbf{Analogy:} Think of sea waves. If one random wave splashes high, you don't panic and think the sea level rose by 5 feet. You want to know the true tide. \\
$x_t$ is the current traffic reading right now. $S_{t-1}$ is the historical smoothed average. \\
$\alpha$ (alpha, e.g., 0.3) controls how much weight we give to the new reading vs. the past. It filters out momentary network spikes so CloudFlip doesn't scale up for a 2-second fluke.
\end{tcolorbox}

\subsection{Formula 2: Discrete Velocity Slope --- The Speedometer}
\begin{equation}
v_t = \frac{S_t - S_{t-k}}{k \cdot \Delta t}
\end{equation}
\begin{tcolorbox}[title=Intuition \& Why We Use It]
\textbf{Analogy:} A car speedometer. It tells you not just your current speed, but how fast you are accelerating. \\
$v_t$ measures the slope (rate of change) of smoothed traffic over a look-back window $k$. If traffic was 8,000 RPM 2 minutes ago and is 12,000 RPM now, the slope is $+2,000$ requests per minute per minute. When slope breaches threshold $\theta_v$, CloudFlip knows a surge is accelerating.
\end{tcolorbox}

\subsection{Formula 3: Lead-Time Horizon Projection --- Looking Down the Road}
\begin{equation}
\hat{x}_{t+H} = S_t + v_t \cdot H
\end{equation}
\begin{tcolorbox}[title=Intuition \& Why We Use It]
\textbf{Analogy:} If you are driving at 60 km/h and accelerating at 10 km/h every minute, you can project where you will be 15 minutes down the road. \\
$H$ is the lead time horizon (15 to 30 minutes). $\hat{x}_{t+H}$ is the predicted traffic $H$ minutes in the future. Because EC2 instances take 5 minutes to boot, ordering capacity 15--30 minutes ahead guarantees the servers are ready before queues ever fill up.
\end{tcolorbox}

\subsection{Formula 4: Target Capacity Fleet Sizing --- Ordering the Right Fleet}
\begin{equation}
\hat{N}_{t+H} = \max\left( N_{\min}, \left\lceil \frac{\hat{x}_{t+H}}{C_{\text{inst}} \cdot \mu_{\text{target}}} \right\rceil \right)
\end{equation}
\begin{tcolorbox}[title=Intuition \& Why We Use It]
$C_{\text{inst}} = 2,500$ RPM (how many requests 1 server can handle safely). \\
$\mu_{\text{target}} = 0.65$ (safety buffer: we aim for 65\% CPU utilization so servers never choke). \\
$N_{\min} = 4$ (baseline floor: the cluster never scales below 4 nodes so cold-starts are prevented). \\
We divide projected demand by per-instance safe capacity, round up to the next integer ($\lceil \dots \rceil$), and ensure it is at least $N_{\min}$.
\end{tcolorbox}

\subsection{Formula 5: Little's Law \& Queueing Latency SLA}
\begin{equation}
L = \lambda W
\end{equation}
\begin{tcolorbox}[title=Intuition \& Why We Use It]
\textbf{Little's Law} is the golden rule of queueing theory: \\
$L$ = average number of requests waiting in queue. \\
$\lambda$ (lambda) = arrival rate of incoming requests. \\
$W$ = average waiting time (latency). \\
If arrival rate $\lambda$ exceeds service capacity, queue length $L$ explodes to infinity, causing wait time $W$ (tail latency) to spike to 4,200~ms! CloudFlip ensures that total cluster capacity $\mu$ expands before $\lambda$ overloads it, keeping $W \le 180\text{ ms}$.
\end{tcolorbox}

\section{The Core Numbers to Know by Heart}
\begin{itemize}[itemsep=2pt]
  \item \textbf{Baseline Fleet:} 4 instances ($N_{\min}=4$), serving 8,000 RPM (2,000 RPM/instance).
  \item \textbf{Surge Profile:} 6$\times$ flash surge from 8,000 to 48,000 RPM within 120 seconds.
  \item \textbf{Instance Capacity:} 2,500 RPM per instance; 20 instances required for 48,000 RPM.
  \item \textbf{Reactive Baseline Result:} Tail latency p99 exploded to \textbf{4,200 ms}; HTTP error rate \textbf{3.8\%}.
  \item \textbf{CloudFlip Result:} Tail latency p99 held within \textbf{180 ms}; HTTP error rate cut to \textbf{0.1\%}.
  \item \textbf{Lead Time:} \textbf{15 to 30 minutes} advance pre-warming.
  \item \textbf{FinOps Waste Audit:} Avoided roughly \textbf{\$670/month} in idle EC2/EBS waste.
\end{itemize}

\section{Five Killer Interview Questions \& How to Answer Them}

\textbf{Q1: Why did you use EMA and derivative slope instead of deep learning (LSTM / Transformers)?} \\
\textit{Answer:} ``In production telemetry streams, deep learning models introduce significant inference latency, high GPU/memory overhead, and require thousands of hours of historical training data that young startups or microservices don't have. CloudFlip's EMA plus velocity slope executes in sub-millisecond CPU time, requires zero offline training, and detects acceleration ramps immediately on streaming data.''

\textbf{Q2: What happens if a surge is a pure step function with no gradual ramp?} \\
\textit{Answer:} ``That is an honest, acknowledged limitation. A slope forecaster requires at least a brief ramp to calculate velocity. If a burst is an instantaneous 10$\times$ jump with zero lead time, CloudFlip falls back to its closed-loop reactive guardrail: the moment p99 latency approaches 180~ms, an emergency step-up trigger immediately scales the fleet. At Princeton, I want to formally prove worst-case bounds for this exact predictive-reactive boundary.''

\textbf{Q3: Why did you submit to IEEE CCGrid 2027 and archive on Zenodo?} \\
\textit{Answer:} ``Zenodo provided immediate, permanent open-access attribution with a DOI (10.5281/zenodo.23136467) without expensive author fees. Following that, I submitted our full manuscript to IEEE CCGrid 2027 under Track 5 (Performance Modeling, Analysis, and Optimization, Paper \#39), which is the premier IEEE/ACM venue for cloud autoscaling and queueing dynamics.''

\textbf{Q4: How does your work connect to Prof. Ravi Netravali's research?} \\
\textit{Answer:} ``In Apparate (SOSP '24), Prof. Netravali's group explored dynamic early exits in ML serving to shed compute during peak load. I want to explore co-designing infrastructure autoscaling with application-level adaptation: when an unexpected burst breaches predictive lead times, how can CloudFlip signal runtime adaptations like Apparate to temporarily reduce per-request compute demand and keep tail latency bounded while new instances warm up?''

\textbf{Q5: How does your work connect to Prof. Amit Levy's research?} \\
\textit{Answer:} ``In Unum (NSDI '23), Prof. Levy demonstrated that replacing centralized serverless orchestrators with decentralized runtime libraries running alongside functions cuts coordination delay. I want to investigate whether CloudFlip's autoscaling policies can be compiled directly into local function runtimes so worker nodes make local queue-draining decisions without waiting for centralized CloudWatch polling loops.''

\end{document}
